\chapter{可解释性}
\label{chap:trust-interpretability}

邮件被模型拦截后，用户希望知道哪些内容触发了判断；开发者则想知道模型是否依赖了不稳定的页脚；如果确实误拦，用户还需要知道如何补救。这三个问题分别涉及单次行为、整体规律和可行行动。给出一张特征热力图，并不自动同时回答它们。

\term{可解释性}（interpretability）关注人能否依据适当的表示和证据理解模型行为。第~\ref{chap:trust-robustness}章通过改变模板发现了可能的捷径，但一次测试失败还不足以定位模型依赖的具体因素。前一章对模型行为作出评估并设置约束，这里进一步追问模型依据什么产生行为，以及哪些证据能够支持这种判断。本章由此确定解释目的，再分别讨论整体与单次行为。解释的质量需要在使用前得到检查；决策补救则进一步引入行动成本、可行性和因果约束。第~\ref{chap:trust-reliability}章检验预测是否可靠，本章研究支持理解和质疑的依据。

\section{模型的可解释性}
\label{sec:interpretability-purpose}

\subsection{解释对象决定证据}

解释至少需要明确对象、使用者和用途。开发者排查错误可能需要内部激活和训练样本；普通使用者判断一次结果可能只需要相关因素及其限制；受影响者寻求补救则需要知道可以采取什么行动。信息更多不一定更有用，解释必须能够支持当前判断。

还应区分模型的行为规律、计算机制和现实因果关系。观察年龄与模型输出相关，可以描述行为；改变模型输入中的年龄字段，可以检查模型是否使用它；在现实中改变年龄通常不是可行行动，更不能据此推断现实结果如何变化。三种论断的对象不同，不能只因为都使用“影响”一词就互相替换。

\subsection{整体解释与单次解释}
\label{sec:interpretability-global-local}

\term{整体解释}（global explanation）是在明确的一组输入或参考分布上，用人能够理解的形式概括模型怎样作出判断。它可以是一套可执行规则、一组特征响应曲线，或关于某类样本的概念作用统计。这里的“整体”限定解释希望覆盖的范围，不表示解释一定包含全部内部计算，也不表示一个平均数已经适用于每个输入。解释邮件模型时，参考范围可以是部署中的全部邮件，也可以是事先声明的某种语言或某类业务邮件。

\term{单次解释}（local explanation）则把目标固定为一次输出，例如这一封邮件为什么被拦截。它可以比较当前输入附近的变化、相对于某个基准的分数差，或者当前运行中某项内部计算的作用。整体与单次解释的共同点是都要把模型行为和可核查的证据联系起来；区别在于待解释的范围和需要支持的论断。整体解释需要说明规律在参考人群中的覆盖程度，单次解释需要说明参照输入、邻域或具体运行条件。

两者能够相互连接。若已有准确的整体加性规则，将当前特征代入，就可以追踪一次预测；若已有许多单次归因，则可以在固定解释定义后统计贡献分布，形成总体摘要。但一次解释不能推出普遍规律，平均归因也可能掩盖方向相反的子群。对模型作输入输出查询，既能得到整体曲线，也能得到单次解释；是否读取内部参数是另一个划分维度，不与整体、局部一一对应。

\begin{example}[同一规则的整体与单次说明]
\label{ex:interpretability-global-local}
设一个教学用邮件评分器读取两个二元特征：$x_1$表示出现支付要求，$x_2$表示出现可执行附件。其规则为$f(x_1,x_2)=1+2x_1+4x_2$，分数不小于6时隔离邮件。分数只是构造的模型打分，不是概率。整体说明需要交代基准分、两项增量和隔离阈值，因此覆盖四种特征组合。对当前邮件$(1,1)$，单次说明则是$1+2+4=7$，超过阈值；相对于$(0,0)$，两项输入分别贡献2和4。这一次代入不说明真实模型在其他组合上也按同样规则计算，整体主张还需要相应结构或查询证据。
\end{example}

\subsection{解释的表示与证据条件}

一个稀疏线性模型或短规则列表通常便于追踪计算，但特征本身仍需有明确含义。若输入是一万个难以解释的嵌入维度，线性结构也未必让使用者理解决策。反过来，复杂模型的某个局部行为可能能够用简单规则准确描述，却不能据此概括整个模型。

解释还会受到访问条件限制。只有输入输出查询时，主要依赖行为近似；能够读取参数和激活时，可以提出机制假设；能够重新训练时，才能更直接地检验某些训练数据影响。后续方法应在这些信息条件下理解，而不是将所有解释视为同一层次的证据。

解释通常还隐含一个比较问题。邮件为什么被拦截，可以指“为什么这一封被拦而相似邮件没有”，也可以指“为什么拦截而不是交给人工复核”。前者比较输入，后者比较行动；它们需要的分数、背景样本和约束不同。即使选择同一种归因算法，改变比较基准也可能改变答案。因此，一份可核查的解释应说明所解释的输出、参照对象、适用输入范围和允许改变的因素。

\section{整体行为的解释}
\label{sec:interpretability-global}

整体行为可以从模型结构、输入和语义三个方面理解。模型结构说明预测怎样由规则、函数或计算分支组成；输入分析说明改变哪些可观察因素会怎样改变输出；语义解释则把词元、像素或特征组合转换成读者关心的概念，考察模型是否利用这些含义。三者可以结合：一个模型可以采用可读结构，以输入特征计算，并通过语义概念表达其中的依据。下面按主要解释对象安排方法，不把这三个方面当成互斥的模型类型。

\subsection{基于模型结构的解释}
\label{sec:interpretability-global-response}

如果要让读者理解“模型遇到不同输入时按照什么规则计算”，可以从一开始就选择可读结构，或者在模型训练后用简单模型近似它。前者的规则属于实际预测器，后者的规则属于解释用的代理。两者都能提供整体说明，但前者需要检验任务表现与结构可理解性，后者还必须检验是否忠实于原模型。

\term{广义加性模型}（generalized additive model, GAM）用若干单特征函数组成预测，使每项特征的作用可以单独查看，例如
\begin{equation}
 f(\bm{x})=\beta_0+\sum_{j=1}^{d}f_j(x_j).
 \label{eq:interpretability-gam}
\end{equation}
这里$\beta_0$为基准项，$f_j$描述第$j$个特征的效应。回归时$f$可以直接作为预测值；二分类时，可令它表示对数几率，再用逻辑函数转换为概率。固定各函数的中心化约定，例如令其在训练分布上的平均值为零，可以分别查看变化趋势。曲线的陡峭区域说明模型对该特征变化敏感，数据稀少区域则需要同时展示样本覆盖，避免把不稳定估计误读为确定规律。

只允许相加会遗漏交互。邮件中的“支付”与“附件”单独出现都可能正常，同时出现才增加风险。Lou等研究的带成对交互的加性模型在单变量项之外增加少量$f_{j,k}(x_j,x_k)$，用二维曲面表达一项特征的作用如何随另一项改变。训练时既要选择哪些交互值得保留，也要控制各项复杂度；解释时则需要同时检查单变量曲线和交互曲面\scite{lou2013ga2m}

规则和树结构同样可以追踪决策路径，但规则数量、重复条件和层次深度会增加理解负担。把一棵树压缩成几条常见路径，也会漏掉低频分支。选择可理解结构时应在当前任务上同时检查预测损失与人能否完成预期判断，不能预先断言简单模型一定足够，也不能未经比较便认定复杂模型不可替代。

构建这类实际预测器时，训练数据提供真实标签，学习算法在加性结构、函数平滑性或树深度等限制下拟合任务。解释阶段固定训练结果，显示各项函数、阈值和交互，并说明如何将它们组合成输出。也就是说，可读的不是训练算法的名称，而是训练后能够被逐项追踪的计算规则。

对于已经训练好的复杂模型，可以另外建立\term{整体代理}（global surrogate），即在声明的参考分布上模仿原模型的简单预测器。具体做法是从该范围收集输入，查询原模型的分数，用这些分数作为拟合目标，再训练浅树、规则表或加性模型。设原模型为$f$，可读模型族为$G$，复杂度惩罚为$\Omega$，$\lambda\geq0$控制拟合误差与复杂度的权衡，则一种构造为
\begin{equation}
 \hat g=\argmin_{g\in G}\left\{
 \frac1n\sum_{i=1}^n\bigl[f(\bm{x}_i)-g(\bm{x}_i)\bigr]^2
 +\lambda\Omega(g)\right\}.
 \label{eq:interpretability-global-surrogate}
\end{equation}
这里使用的是原模型输出$f(\bm{x}_i)$，而非真实标签$y_i$；对真实标签重新训练一个小模型，不能自动解释原模型。分类时也可以拟合概率向量或类别，但选择哪一种目标，会决定代理保留分数变化还是仅保留决策边界。

得到代理后，需要展示其实际规则，并在另外保留的查询样本上比较$g$与$f$。例如，浅树如果先检查附件，再检查支付要求，就能在示例~\ref{ex:interpretability-global-local}的四种组合上复述隔离规则；只写“有附件就隔离”则在$(0,1)$上与原模型不符。平均一致率很高而遗漏少数业务分支时，应保留分组误差，或将解释范围缩小到真正拟合良好的区域。增加代理复杂度通常能减少近似误差，却可能使解释重新变得难读。

\subsection{基于输入的解释}
\label{sec:interpretability-global-dependence}

整套代理试图给出一个能够重新计算预测的模型；如果只想知道“哪些输入怎样影响判断”，可以直接从原模型提取作用摘要，无须先学出整套规则。这就是本小节与上一小节的关系。摘要主要有两个互补问题：特征改变时，分数怎样响应；破坏该特征信息后，任务表现损失多少。前者描述作用形状和方向，后者度量固定模型对信息的性能依赖。

\paragraph{响应形状与交互}

\term{部分依赖}（partial dependence, PD）分析考察把某个特征固定后，对其余特征取平均的输出：
\[
 \operatorname{PD}_j(t)
 =\mathbb E_{\bm{x}_{-j}}f(t,\bm{x}_{-j}).
\]
其中$\bm{x}_{-j}$表示除第$j$个特征外的输入，$t$为要考察的特征取值。

实际构建部分依赖曲线时，选取一组特征取值$t$，把参考样本的第$j$项都替换为$t$，逐一查询模型，再求平均：
\[
 \widehat{\operatorname{PD}}_j(t)
 =\frac1n\sum_{i=1}^n f(t,\bm{x}_{i,-j}).
\]
将这些均值随$t$画出，就得到上升、下降、阈值或饱和的形状。若同时考察两个特征，则对二维取值网格作同样替换和平均，以检查一个特征的作用是否随另一个改变。只展示一条平均曲线时，还应检查各输入的响应曲线是否方向一致。

沿用示例~\ref{ex:interpretability-global-local}，若四种输入组合在参考群体中等概率出现，则$\operatorname{PD}_1(t)=3+2t$，$\operatorname{PD}_2(t)=2+4t$。这些曲线概括支付要求和附件对平均分数的作用，却没有给出每封邮件的分数，更没有使用真实标签判断模型是否正确。若参考群体改变，平均曲线的截距也可能改变；在这个没有交互的例子中，斜率仍由规则中的2与4决定。

这种替换也可能组合出现实中不存在的样本。若所有观测满足$x_1=x_2$，模型为$f(x_1,x_2)=x_1-x_2$，它在观测数据上恒为零；单独改变$x_1$却会得到上升趋势。趋势描述的是模型对这些组合的响应，不能直接解释为观测人群中的实际效应。

条件平均使用给定特征后的其余特征分布，避免部分不现实组合，却会混合特征关联。因此，即使画出的都是一条“特征—输出”曲线，也需要说明固定了什么、对什么分布取平均。

为了减少跨越数据支持区域的外推，Apley与Zhu提出\term{累积局部效应}（accumulated local effects, ALE）：只在某特征当前取值附近估计小幅变化的平均作用，再将这些局部作用沿该特征累积。对适当光滑的模型，一阶形式为
\begin{equation}
 \operatorname{ALE}_j(t)
 =\int_{t_0}^{t}
 \mathbb E\!\left[
 \frac{\partial f(s,\bm{x}_{-j})}{\partial s}
 \,\middle|\,x_j=s\right]\mathrm ds-C_j,
 \label{eq:interpretability-ale}
\end{equation}
其中$t_0$为积分起点，$C_j$用于中心化。实际计算通常把特征值分箱，在每箱已有样本上比较箱边界处的输出，再累加平均差值；不必求模型导数\scite{apley2020ale}

部分依赖平均的是整个预测值，累积局部效应平均的是局部变化，这正是两者的差别。后者降低了不现实组合造成的外推问题，但分箱太粗仍会跨过无数据区域，分箱太细又可能估计不稳。它也不是从观测数据识别现实因果作用的方法：同一模型函数在数据流形以外如何延伸，仍可能影响其局部偏导。若要理解个体差异，可以保留每个样本随特征变化的曲线，检查总体平均掩盖了哪些交互。

\paragraph{任务表现对信息的依赖}

响应曲线说明分数如何变化，却没有使用真实结果来判断这些变化是否有助于完成任务。若要检查模型是否依赖邮件页脚来正确分类，需要比较保留与破坏这项信息时的预测损失。这个问题以任务表现为参照，因此不同于读取曲线斜率或系数大小。

置换重要性固定已训练的模型，将测试集中某项特征在样本间打乱，再比较平均损失与原始损失。实际检验应在同一批测试样本上作成对比较，重复置换，报告损失增量的均值和波动。若打乱页脚后错误明显增加，得到的是模型在这种信息破坏下的性能依赖证据；它既不证明页脚具有现实因果作用，也不直接说明这种依赖合理。

把置换过程写开，可以更清楚地看到解释的对象。设$\pi$为样本索引的随机排列，$\bm{x}_i^{(j,\pi)}$只把样本$i$的第$j$项替换成$x_{\pi(i),j}$，则一次损失增量为
\begin{equation}
 \widehat I_j^{(\pi)}
 =\frac1n\sum_{i=1}^n\left[
 \ell\bigl(f(\bm{x}_i^{(j,\pi)}),y_i\bigr)
 -\ell\bigl(f(\bm{x}_i),y_i\bigr)\right].
 \label{eq:interpretability-permutation}
\end{equation}
需要真实标签的是损失比较，原模型在置换前后始终不重训练。以零一损失为例，若一个教学测试集有100封邮件，原来误判10封，一次置换附件字段后误判25封，则这次增量为0.15。重复置换得到的均值和波动，分别描述该破坏规则下的平均性能变化与置换随机性；对新邮件总体的推广还要考虑测试样本的不确定性。

关联特征使参照尤其重要。若页脚与发送方始终配套，单独打乱页脚可能制造不现实组合，损失增量同时包含这种分布外效应。按发送方分组后置换，考察的是给定发送方后页脚剩余信息的作用；把相关特征作为一组置换，则考察这组信息的共同作用。若模型只读取两个相互替代的特征之一，另一项的置换重要性可以为零，但不能由此断言它在数据中不含预测信息。这里检验的是固定模型的使用方式，删除特征后重新训练比较的是学习系统能否找到替代依据，两者需要分别报告。

曲线与重要性因而不应互相代替。分数对某特征变化明显，但变化没有跨过分类阈值时，零一损失的重要性可能很小；模型若错误利用某项信息，破坏它甚至可能降低损失。响应曲线回答模型如何变化，损失增量回答这种依赖对所选任务指标意味着什么。要判断依赖是否合理，还需结合任务语义和使用要求。

\subsection{基于语义的解释}
\label{sec:interpretability-global-concepts}

输入层的分析仍可能离读者关心的含义很远。例如，“催促付款”可以由许多不同词语表达，同一个词也可能出现在引用或警示中。基于语义的解释用“来源可验证”“要求立即转账”等概念描述依据，再检查这些概念怎样影响输出。构建途径有两种：让模型直接通过命名概念作出预测，或者在已有模型中寻找与概念对应的表示并检验其作用。

Koh等研究的\term{概念瓶颈模型}（concept bottleneck model）先将输入映射为人定义的概念，再只通过这些概念预测结果。它把可查看、可修改的语义接口放进模型的实际计算路径，而不只在预测之后附加描述\scite{koh2020conceptbottleneck}

设概念向量为$\bm{c}\in\mathbb R^m$，概念预测器为$g$，决策函数为$h$，则$f(\bm{x})=h(g(\bm{x}))$。有概念标注时，可以先用概念损失训练$g$，再用预测概念或真实概念训练$h$；也可以联合优化任务损失与概念损失。不同方式改变了任务信号能否反过来调整概念表示，以及训练和使用时概念误差是否匹配。若人工确认邮件来源有效，便可替换相应概念值，观察最终判断是否得到纠正。

要据此给出整套决策规则，还需说明决策函数$h$怎样组合概念，例如展示其加性项、阈值或规则路径。若$h$本身仍是难以理解的复杂网络，概念瓶颈提供的是可检查的语义接口，并没有自动解释整套计算。对概念作用作整体判断，还需在覆盖目标任务的样本上检验，而不能只展示一次纠错。

接口具有名称也不代表语义已经得到保证。概念预测错误会向后传播；遗漏重要概念会限制任务能力；连续概念分数还可能携带超出标注名称的信息。若模型另有绕过瓶颈的输入通路，修改概念后也未必能控制全部依据。验证时应检查概念预测质量、概念替换效果和任务损失，而不能只展示若干容易命名的中间节点。

若不希望重新训练概念瓶颈模型，也可以在已有表示中检验一个明确的概念假设。Kim等提出的\term{概念激活向量检验}（testing with concept activation vectors, TCAV）把概念表示成内部空间中的方向，再检查模型对该方向是否敏感。方法使用带有目标概念和对照概念的样本，在某层激活空间训练线性分隔器，得到指向目标概念一侧的方向$\bm{v}_C$。它再考察类别分数沿该方向的导数，统计一类样本中导数为正的比例。这样，“模型是否使用某种纹理”被转化为可以重复测量的局部敏感性问题\scite{kim2018tcav}

若某层激活为$\bm{h}(\bm{x})$、该层之后的类别分数为$s_k(\bm{h})$，被统计的量是$\nabla_{\bm h}s_k(\bm{h}(\bm{x}))^\top\bm{v}_C$的符号。结果依赖概念样本、对照集和网络层，因此需要更换对照、重复拟合并报告不确定性。概念方向能分开两组样本，可能源自与概念共现的背景；导数为正也只是表示空间中的敏感性，不能直接解释为现实中增加该概念就会提高该类概率。与概念瓶颈相比，TCAV增加了事后检验能力，却没有强制模型只能通过被命名的概念作决策。

具体实施时，需要三组具有不同职责的数据：概念实例及其对照用于学习方向，目标类别或任务样本用于计算作用，另留的概念和任务样本用于复核。固定模型与层之后，记录概念样本的激活，用线性分类器分开概念与对照，并将法向量朝概念一侧定向。随后在每个目标样本上计算类别分数沿该方向的导数。对目标样本集合$D_k$，一种摘要是
\[
 \widehat t_{C,k}
 =\frac1{|D_k|}\sum_{\bm{x}\in D_k}
 \mathbf1\!\left\{
 \nabla_{\bm h}s_k(\bm h(\bm{x}))^\top\bm v_C>0
 \right\}.
\]
若100个待解释样本中有80个方向导数为正，统计量就是0.8，含义是这些样本中有八成在该表示方向上呈正局部敏感性。它不是“该概念贡献了80\%的预测”，也不代表八成预测由它单独决定。更换概念实例和随机对照后仍能重复发现这种差别，才更支持所提出的概念假设。

模型结构、输入和语义的关系也由此明确：结构交代怎样组合证据，输入分析检验具体字段或区域的作用，语义解释说明这些依据对应什么含义。概念瓶颈通过改变结构提供语义接口，概念激活检验则借助已有表示分析语义；同一种方法可以连接多个方面。无论从哪一方面形成整体结论，都要说明样本覆盖、汇总方式以及不同输入间是否存在相反作用。

\section{单次行为的解释}
\label{sec:interpretability-local}

同一次预测可以从三个方面解释。基于推理输入的解释，研究当前输入及其变化与输出有什么关系；基于推理计算过程的解释，研究模型通过哪些计算形成这次输出；基于训练数据的解释，则把当前判断追溯到训练阶段，研究哪些训练经验影响了它。以误拦邮件为例，三个方面分别检查触发判断的邮件内容、处理这些内容的计算环节，以及使模型学会这种判断的训练样本。

相似案例、反事实比较、局部代理和特征归因都属于第一个方面：它们固定已经训练好的模型，通过参照案例、输入变化或贡献分配理解当前结果。推理计算过程的解释进一步检查预测过程中的作用位置；训练数据解释则研究改变学习经历后，同一输入的预测会怎样变化。这里按解释所要说明的对象分组，而不是按算法名称或是否读取内部参数分组。三个方面可以独立提供证据，也可以相互核验。

选择方法前应固定要解释的输出。例如，垃圾邮件分数、最终隔离动作和“隔离而非人工复核”的选择并不相同。若最终动作还受规则阈值或权限约束控制，仅解释分类分数会遗漏后续决策。下面主要把$f(\bm{x})$视为待解释的标量分数，用$\hat y(\bm{x})$表示结合决策规则后得到的类别或动作；多类别模型还需固定所解释的目标类别。

\subsection{基于推理输入的解释}
\label{sec:interpretability-local-input}
\label{sec:interpretability-local-examples}
\label{sec:interpretability-local-counterfactual}
\label{sec:interpretability-local-neighborhood}
\label{sec:interpretability-local-attribution}

这一类解释固定模型，围绕当前输入寻找参照或构造变化，再将输入差异与输出联系起来。相近输入得到相同结果，可以提供共同特征的线索；相近输入得到不同结果，可以提示判断的边界；对多种输入变化作系统比较，则可以刻画响应规律或分配特征贡献。这些比较的共同对象都是推理时的输入与输出，差别在于希望得到案例对照、响应规则还是贡献数值。

并非每种方法都要实际修改原样本：可以检索已有案例，也可以利用模型导数计算局部变化。输入不同而最终类别相同，分数仍可能明显改变；因此，解释还要区分所比较的是类别、分数还是最终动作。

\paragraph{相似案例的对照}

一种直观解释是展示与当前输入相似、并被模型作出相同判断的案例。例如，当前邮件含有支付要求与可执行附件，并被模型隔离，可以展示内容相近、同样包含这两个因素的已隔离邮件，并标明共同片段。读者由此看到当前案例被放到了怎样的一组案例中，而不必先读一组抽象权重。这通常称为\term{基于案例的解释}（example-based explanation）。

构造时，需要先规定相似性：按词语、业务字段还是模型表示比较，哪些差异可以忽略；再从参考库中选出模型结果相同的近邻，呈现共同点、差异和距离。若参考案例有真实标签，还应同时显示标签与预测，避免把一组相似误判展示成判断正确的证据。可以选多个近邻，或挑选有代表性的案例作为原型，减少只展示某个偶然例子的影响。

这种比较首先提供的是理解线索。如果模型本来就按近邻或原型相似度预测，显示实际参与投票的案例及其权重，就能追踪决策的一部分；如果模型独立作出预测后才检索相似案例，检索结果只是事后对照，还不能证明模型用了这些共同特征。需要进一步改变共同特征、比较输出，才能检验这条解释。

Chen等提出的原型部件网络ProtoPNet把案例比较纳入图像分类：比较当前图像局部与学得原型的表示相似度，再加权形成类别分数。原型与训练图像片段关联，使解释能够同时展示“当前哪一块像哪个典型片段”及其对类别分数的作用。这提供了案例比较属于实际预测过程的一种实现\scite{chen2019protopnet}

相似性也需要检查。模型表示可能把背景当作相似依据，人却以为它比较的是物体；邮件模板相似，也不表示业务意图相同。因此应显示相似的具体位置，并核对贡献权重，不能只给两个完整案例便让读者自行推断比较依据。对当前结果的理解，还需要观察相近但结果不同的案例。

\paragraph{不同结果的反事实比较}

相似案例说明“这一输入像哪些同结果案例”，不同结果的比较则追问“哪些差别足以使判断不同”。可以从真实数据中寻找与当前输入接近、但预测不同的案例，逐项比较字段；也可以主动构造新的输入，只改变少量条件并再次查询模型。后者更便于控制其他因素，明确结果依赖的条件。

Wachter等讨论的\term{反事实解释}（counterfactual explanation）寻找接近当前输入、却得到目标输出的另一输入。常见计算形式是最小化距离与目标输出损失的加权和；也可以把达到目标写成约束。它把“为何是这个结果”转化为“哪些条件改变后便不再是这个结果”\scite{wachter2018counterfactual}

设$Y_{\mathrm{target}}$是与当前结果不同的目标结果集合，$d$衡量输入差异，$\bm X_{\mathrm{valid}}$规定本次解释中允许考虑的有效输入，则一种构造为
\begin{equation}
 \min_{\bm{x}'\in\bm X_{\mathrm{valid}}}
 d(\bm{x}',\bm{x}),
 \qquad \hat y(\bm{x}')\in Y_{\mathrm{target}}.
 \label{eq:interpretability-counterfactual}
\end{equation}
实际搜索可以从已有的不同结果案例开始，也可以在允许字段上枚举、使用梯度或优化求解。得到候选后必须重新运行完整预测规则，检查目标结果确实改变，并列出修改了哪些字段。距离可以偏好只改变少数因素；连续变量的尺度、离散变量的取值和字段间的有效组合需要事先说明，否则“最近”可能仅由单位选择决定。有限搜索找到一个候选，不等于已经证明它最小或唯一。

沿用示例~\ref{ex:interpretability-global-local}，当前$(1,1)$得7分而被隔离。保留支付要求、把可执行附件改为非可执行形式，输入成为$(1,0)$，得3分；保留附件、去掉支付要求，输入成为$(0,1)$，得5分，两者都会被放行。这两个反事实分别表明：在另一条件保持时，改变任一项都足以越过当前决策边界。若只展示第一个，就容易让人误以为附件是唯一决定因素。

反事实因此可以定位本次结果所依赖的条件，但结论相对于模型、允许改动和参照成立。真实的不同结果近邻若同时改变多个字段，只能提示差异；受控地改变字段并重算模型，才进一步检验这些改动对预测的作用。即使找到了这样的条件，也尚未证明现实结果由它造成，更没有说明使用者能够实施该改动。前者需要额外的因果依据，后者由第~\ref{sec:interpretability-recourse}节讨论。

案例对照和反事实给出具体参照输入。若希望更系统地说明哪些词语、像素或字段支持当前输出，可以在附近构造多种输入变化，观察输出规律，或者把当前分数相对于基准的差异分配给各项输入。下面的方法共享输入与输出这一比较对象，区别在于拟合附近的响应规则、测量单项变化，还是分配相对于基准的贡献。

\paragraph{附近输入的响应规则}

局部代理方法在一个输入附近查询模型，用简单模型近似这一区域的行为。\term{局部可解释模型无关解释}（local interpretable model-agnostic explanations, LIME）根据与待解释输入的接近程度加权拟合，并限制代理复杂度。邻域如何采样、距离如何定义以及代理误差多大，都会影响解释\scite{ribeiro2016lime}

以邮件为例，可以用词语是否保留构造二元向量$\bm{z}$，通过映射$r(\bm{z})$得到可查询的文本，再拟合稀疏线性代理$g$：
\begin{equation}
 \hat g=\argmin_{g\in G}\left\{
 \sum_{i=1}^{n}w_{\bm x}(\bm{z}_i)
 \bigl[f(r(\bm{z}_i))-g(\bm{z}_i)\bigr]^2
 +\lambda\Omega(g)\right\}.
 \label{eq:interpretability-lime}
\end{equation}
这里$G$为候选代理模型族，$w_{\bm x}$强调与当前邮件接近的扰动，$\Omega$惩罚代理复杂度，$\lambda\geq0$控制惩罚强度。权重范围越窄，代理越侧重当前点附近，却可能缺少足够变化来估计系数；范围越宽，则更难用简单模型覆盖复杂行为。删词可能破坏语法甚至改变邮件目的，因此必须检查生成邻域是否符合解释任务，并在另采样的邻域中报告拟合误差。稀疏系数只是这个加权近似问题的答案。

实际计算时，LIME先把当前邮件表示为可读的词项保留向量，采样若干删词组合，取得原模型对这些组合的分数；再按接近程度加权，用少量非零系数拟合这些分数。最后交付的是所选词项、系数方向和大小，以及代理在声明邻域中的误差。原模型不需要因此重新训练。

若待解释模型恰好是示例~\ref{ex:interpretability-global-local}的加性规则，邻域又包含足以识别两项作用的组合，未加稀疏惩罚的线性拟合便能恢复系数2和4。但真实模型存在交互时，同一词在另一封邮件附近可能得到不同系数。这也说明整体代理与LIME共享“查询后拟合”的思想，差别在于拟合数据和误差权重覆盖整个参考范围，还是围绕当前输入集中。

邻域规则可以告诉人哪些附近输入会改变判断，却未必说明当前输出由哪些证据支持。证据归因针对当前输出标记或分配各项输入信息的作用：有些方法只标记敏感位置，有些方法还要求各项贡献加起来等于某个分数差。后者需要固定比较基准或特征组合的评价规则。归因数值的含义由这些约定决定，不能直接解释为现实原因。

\paragraph{扰动比较与局部敏感性}

只需查询模型的一种做法是遮挡输入：把当前第$j$项替换成规定的参照值，计算原分数减去替换后的分数。对文本可以移除或替换词段，对图像可以替换区域。重复这一操作就得到输入位置的作用图。这里得到的是单项移除效果；有交互或冗余时，所有移除效果之和未必等于总分数差。选择什么替换值，仍是解释定义的一部分。

对于可微分分数$f$，梯度说明小幅改变输入时输出的局部敏感性。它可能在饱和区域接近零，即使该特征曾对输出形成过程十分重要。仅有终点敏感性还不足以分配当前分数相对于基准的变化，需要另外规定沿途如何累计。

\paragraph{相对于基准的贡献分配}

\term{积分梯度}（integrated gradients, IG）沿基准$\bm{x}_0$到输入$\bm{x}$的路径累计导数：
\begin{equation}
 \operatorname{IG}_j(\bm{x})
 =(x_j-x_{0,j})\int_0^1
 \frac{\partial f(\bm{x}_0+t(\bm{x}-\bm{x}_0))}
      {\partial x_j}\,\mathrm dt.
 \label{eq:interpretability-ig}
\end{equation}
在沿路径满足相应微积分条件时，各维贡献之和等于$f(\bm{x})-f(\bm{x}_0)$。这个等式说明贡献解释了相对于基准的分数差，却不赋予基准或路径现实因果含义。Sundararajan等从归因公理出发提出这一方法；黑色图像、空文本或平均输入是否适合作为基准，仍需要由任务判断\scite{sundararajan2017ig}

这一求和关系可由链式法则理解：沿直线路径对$f$求导，得到各维偏导与位移的乘积之和；对路径参数积分后，左端就是两端分数差。因此，积分梯度能够累计路径中经过的敏感区域，而最终点的梯度只描述终点。对于离散文本，在词嵌入之间插值只是计算路径，路径上的向量未必对应任何合法句子。数值积分的步数也应通过求和残差检查，不能把数值误差当成模型遗漏的证据。

实际求积分梯度时，选择基准与当前输入之间的若干插值点，在每点计算目标分数的输入梯度，求平均后乘以各维输入差。示例~\ref{ex:interpretability-global-local}在$(0,0)$到$(1,1)$的路径上梯度始终为$(2,4)$，因此贡献为2和4，总和为$7-1=6$。这说明了如何从模型计算得到解释数值；换成另一个基准后，被分配的分数差也会改变。

Shapley值将贡献分配理解为不同特征组合带来的增量。若$N=\{1,\ldots,d\}$，$v(S)$表示保留特征子集$S$时的约定价值，则第$j$个特征的分配为
\begin{equation}
 \phi_j=\sum_{S\subseteq N\setminus\{j\}}
 \frac{|S|!(d-|S|-1)!}{d!}
 \bigl[v(S\cup\{j\})-v(S)\bigr].
 \label{eq:interpretability-shapley}
\end{equation}
SHAP（Shapley additive explanations）利用这类加性归因解释模型预测，即把相对于基准的分数变化分配给各项特征。关键仍是$v(S)$的定义：缺失特征如何填补、是否使用条件分布，会改变被解释的问题和结果\scite{lundberg2017shap}

一种选择是将未保留特征从背景分布采样，令$v(S)=\mathbb E[f(\bm{x}_S,\bm{X}_{-S})]$；另一种选择按已观测特征条件采样。前者更直接研究模型在人工组合上的响应，后者保留数据关联，但可能把贡献分给模型未直接读取、却能预测其他特征的变量。两者回答不同问题，并不存在脱离任务的唯一“真实贡献”。精确计算还要处理指数数量的子集，实际算法需要利用模型结构或采样近似，报告背景集与近似误差。

可以把Shapley计算理解为反复随机排列特征的进入顺序。在每一种顺序中，从背景开始逐项加入当前输入的特征值，记录加入第$j$项前后的价值差；对所有顺序取平均，就得到式~\eqref{eq:interpretability-shapley}。特征很多时采样一部分顺序近似。其产物仍是一组当前输入的贡献，区别于积分梯度的是比较不同特征组合，而不是沿一条连续输入路径累计梯度。

\begin{example}[共同信息的贡献分配]
\label{ex:interpretability-redundancy}
在一个构造的两特征问题中，令$v(\varnothing)=0$，$v(\{1\})=v(\{2\})=v(\{1,2\})=1$。两项信息互相替代，任一项都足以达到价值1。Shapley分配对每项给出$1/2$：它平均了该特征先进入或后进入的增量。这个数不表示删除任意一个特征会让完整模型输出降低一半。
\end{example}

\paragraph{支持区域的定位}

图像中的归因还可以以区域而非输入坐标为单位。此时人常常只需要先找到支持当前类别、值得进一步检查的位置。Grad-CAM将目标类别分数对某层卷积特征图的梯度作空间平均，以此给不同通道加权，再保留加权和的正值区域。若第$k$个通道的特征图为$\bm A^k$，空间位置数为$Z$，目标类别分数为$s_c$，则
\begin{equation}
 \begin{aligned}
 \alpha_k^c&=\frac{1}{Z}\sum_{i,j}
 \frac{\partial s_c}{\partial A^k_{i,j}},\\
 L^c_{i,j}&=\max\!\left(0,\sum_k\alpha_k^c A^k_{i,j}\right).
 \end{aligned}
 \label{eq:interpretability-gradcam}
\end{equation}
梯度指出哪些通道对当前类别分数敏感，特征图提供这些通道在空间中的响应，二者结合形成粗粒度定位。Selvaraju等用这一机制解释多种卷积网络的类别相关区域\scite{selvaraju2017gradcam}

定位图与积分梯度的目标不同：前者强调某层表示中支持目标类别的区域，后者分配相对于输入基准的分数差。Grad-CAM的分辨率受所选层限制，空间平均会丢失通道内的差异，保留正值还会隐藏反对该类别的证据。上采样后的清晰边缘不能被理解为像素级因果结论。同一图像应检查不同目标类别和对照区域，避免把“看起来像目标物体”当成唯一依据。

\subsection{基于推理计算过程的解释}
\label{sec:interpretability-local-mechanism}

前一小节把当前结果与推理输入联系起来，但尚未说明这些信息在模型中怎样参与计算。推理计算过程的解释固定已训练的参数，检查预测过程中哪些表示、组件或通路影响了输出。读取内部激活后，可以寻找编码某种概念的方向，再通过替换、屏蔽或恢复激活检查行为变化。能够从激活解码某个信息，说明该信息存在于表示中，却不说明当前决策使用了它。干预更接近检验计算作用，但也可能制造训练时从未出现的内部状态，因此需要合适的对照。

一种具体做法是保留原输入的干净运行，再构造使目标行为受损的输入，把干净运行中某个位置的激活替换进受损运行。如果目标输出恢复，便获得该位置在这一干预设置中参与行为的证据。Meng等的因果追踪研究用这种思路定位语言模型事实关联的相关计算，并进一步以参数编辑检验机制假设。局部恢复不能证明信息只储存在该位置，因为替换值可能携带多个因素，其他通路也可能冗余地保存信息\scite{meng2022rome}

把这一方法落实到一次预测，需要分别保存干净运行的激活和目标分数，进行受损运行，再只在选定层、位置或组件上换回干净激活，继续向前计算。设三次目标分数为$s_{\mathrm{clean}}$、$s_{\mathrm{corrupt}}$和$s_{\mathrm{patch}}$，且前者大于受损分数，可以用
\begin{equation}
 r_{\mathrm{restore}}
 =\frac{s_{\mathrm{patch}}-s_{\mathrm{corrupt}}}
        {s_{\mathrm{clean}}-s_{\mathrm{corrupt}}}
 \label{eq:interpretability-restoration}
\end{equation}
描述此次替换恢复了多少分数差。在一个构造例子中，三个分数分别为4、0和3，恢复比例就是0.75。它表示当前替换实验恢复了四分之三的差值，不表示该组件存储了四分之三的知识。比例也可能为负或超过1，分母太小时则很不稳定。逐个改变位置形成恢复图，再用其他输入配对和对照替换复查，才能提出较有依据的计算通路假设。

这与输入归因的区别在于解释的论断：归因把输出差分配给输入因素，机制干预询问内部哪一段计算在规定操作下影响输出。二者都可能使用梯度或激活，因此不能仅凭“使用了内部信息”来划界。例如Grad-CAM虽然读取内部特征图，交付的主要是输入区域定位；激活替换则直接比较改变内部计算前后的行为。输入归因有助于提出机制假设，内部干预再检验其中一部分，两者没有自动等价关系。

干预需要把必要性与充分性分开。移除某组件使行为消失，支持它在当前设置中的必要性；仅恢复该组件便挽救行为，支持它在特定背景下的恢复作用。改变污染输入、替换位置或评分目标可能改变结论。通过多种输入配对、等强度随机替换和组合消融检验，可以逐步排除简单的替代解释；单个成功案例尚不足以得到完整计算图。

对一次工具调用或一段生成过程，解释单位也可以扩展为行动序列。此时应追踪哪一步引入了关键证据，哪一步改变了后续选择，以及恢复某一步是否能够纠正结果。只看最终文本可能遗漏真正造成错误的中间行动。

\subsection{基于训练数据的解释}
\label{sec:interpretability-local-training}

前两小节固定已经训练好的模型，分别解释推理输入的作用和内部计算。若怀疑某次误拦源于训练集中的错标样本，就需要把当前结果追溯到学习阶段：改变哪些训练样本后，模型对这封邮件的判断会发生变化？训练样本同样包含输入，但它在这里通过训练过程影响模型参数；在监督学习中，样本标签也参与这种影响。分析时保持待解释邮件不变，改变训练样本、标签或权重，研究学得模型的预测变化。因此，训练数据影响不同于前文对当前邮件删词后直接重算分数。

删除样本后重新训练是直接参照，但代价较大。影响函数在局部光滑条件下，用最优解对样本权重的导数近似这种变化。设训练目标为$\hat R(\bm{\theta})$，局部最优解为$\hat{\bm{\theta}}$，该点处可逆的海森矩阵（Hessian matrix）为$\bm H$。将目标改为$\hat R(\bm{\theta})+\epsilon\ell(z,\bm{\theta})$，并沿原局部解附近的平稳点分支得到$\hat{\bm{\theta}}_\epsilon$，则
\[
 \left.\frac{\mathrm d\hat{\bm{\theta}}_\epsilon}
 {\mathrm d\epsilon}\right|_{\epsilon=0}
 =-\bm H^{-1}\nabla_{\bm{\theta}}
       \ell(z,\hat{\bm{\theta}}).
\]
这个关系来自对一阶最优条件求导。它依赖局部解和光滑性，大幅删除、非唯一解或重新训练路径变化可能使近似失准。Koh与Liang研究了这种近似在模型预测追溯中的应用；训练数据归因仍应通过部分重训练或受控修改验证，而不能把近似分数直接当作数据来源的确定证明\scite{koh2017influence}

若要解释测试样本$z_*$的损失，还需左乘$\nabla_{\bm\theta}\ell(z_*,\hat{\bm\theta})^\top$。由此得到的负梯度内积经过逆海森矩阵调整，衡量增加训练样本权重时测试损失的局部变化。它不同于找输入最相似的训练样本：外观相近的两封邮件可能产生不同梯度，而一个不相似的错标样本也可能显著改变决策边界。大型模型通常通过海森矩阵与向量的乘积近似求解线性系统，不必显式存储整个逆矩阵；但近似计算并不会消除局部最优假设。

若训练没有达到稳定最优点，可以改为追踪训练过程。TracIn考察使用某训练样本更新参数时，待解释样本的损失改变多少。对单样本梯度下降，一阶展开给出
\[
 \ell(z_*,\bm\theta^{(t)})-\ell(z_*,\bm\theta^{(t+1)})
 \approx\eta_t
 \nabla\ell(z_*,\bm\theta^{(t)})^\top
 \nabla\ell(z_i,\bm\theta^{(t)}).
\]
其中$\eta_t$为第$t$步学习率，梯度均对模型参数求取。Pruthi等据此在保存的训练检查点上累积梯度内积，近似识别有助于降低或增加该损失的训练样本。它避开了最终海森矩阵求逆，但需要检查点，并依赖更新步长、批处理和优化器；检查点近似也不等于删除样本后从头训练的反事实结果\scite{pruthi2020tracin}

\section{解释的质量}
\label{sec:interpretability-quality}

相似案例、反事实、输入贡献和计算追踪都给出了理解结果的线索，还需检验这些线索能否支持所作的说明。本章按五个问题组织评价：解释是否符合被解释模型，解释覆盖了声称的行为范围，重复计算与合理变化下是否稳定，人能否正确理解，以及这种理解是否改善实际使用。它们分别对应忠实性、完整性与覆盖、稳定性、可理解性和使用有效性。这里是一组需要分别核对的维度，不存在一个能够无条件汇总所有解释的通用分数。

评价应先固定模型版本、目标输出、输入范围、基准或干预规则和解释使用者。整体规则、特征摘要、单次归因与内部机制的论断不同，忠实性测试也必须匹配；对所有方法一律画一张删除曲线，会检验错对象。下面先说明每个维度的判断含义，再给出可以实际执行的比较。

\subsection{忠实性}
\label{sec:interpretability-quality-fidelity}

\term{忠实性}（fidelity）指解释是否准确反映它声称描述的模型行为。被解释模型在任务上可能错误，解释却可以忠实地暴露这种错误依据；因此，解释忠实性与模型预测正确率需要分开。

对相似案例，应先复核相似度和输出是否符合检索条件，再区分案例只是事后参照，还是实际参与了预测。近邻模型可以重新计算邻居投票，原型模型可以核对相似度与类别权重；事后检索若声称共同特征解释了当前结果，还需通过改变这些特征检验输出。只证明两例相似，不足以检验关于模型依据的主张。

对反事实，直接检查候选经过完整系统后是否得到目标结果，并报告输入距离、修改数量和有效性约束。若声称某项改动不可缺少，可以撤销该改动后重算；多项改动存在交互时，需要比较其组合。还应搜索其他相近反事实，检查解释是否只选了一条有利路径。是否能通过现实行动实现，是补救有效性的额外检验，不能用预测翻转代替。

对整体或局部代理，在未参与拟合的查询样本上同时运行原模型$f$与代理$g$，比较分数误差或决策一致率。取非负样本权重$w_i$，标量分数的加权均方误差可以写为
\[
 \widehat E_{\mathrm{fid}}
 =\frac{\sum_i w_i\,[f(\bm{x}_i)-g(\bm{x}_i)]^2}
        {\sum_i w_i},\qquad \sum_iw_i>0.
\]
整体评价的样本与权重对应声明的参考范围，局部评价则对应当前输入的邻域。若解释只声称复述类别，就报告预测类别一致的比例；若声称复述概率，类别一致仍不足以检验分数误差。按重要业务类型分别计算，可以发现总体误差掩盖的分支。

对特征归因，要将“贡献分配是否满足其定义”和“是否预测干预效果”分开。积分梯度的贡献和应接近基准分数差，SHAP应符合声明的背景和特征缺失规则；这些是计算一致性检查。若进一步主张高归因输入对输出更关键，可以按归因排序逐步删除高分部分，记录目标分数随删除比例的变化，并与随机删除、低分部分先删除比较。反向实验从基准开始逐步插入高分部分，检查分数怎样恢复。删除、插入曲线及其曲线下面积给出相应扰动规则下的证据；并非所有满足分配公理的方法都应在任意删除规则下取得最优结果。

这类检验必须控制输入有效性和改动量。删除文字可能改变语义，遮挡图像可能产生异常块，需要与同样大小、同样替换方式的对照比较。既含支持也含反对证据时，应分别检查正负作用；不能把移除反对证据后分数升高误写成解释失败。对于部分依赖或概念作用摘要，可以在新的参考样本上重算曲线或概念敏感性，并检查重要差异是否复现。

Adebayo等提出的显著性图健全性检验，直接检查解释是否依赖它声称解释的模型与数据。一类检验逐步随机化网络参数，观察解释是否改变；另一类在随机打乱标签后训练网络，再比较解释。如果一幅图在预测机制已经改变后仍几乎不变，它可能主要反映输入边缘或架构先验，不足以解释模型学到的输入输出关系。研究报告的失败针对被测试方法与设置，不能据此宣称所有显著性方法都无效\scite{adebayo2018sanity}

通过随机化检验只是排除一种明显错误。机制解释还应在新的输入配对上预测屏蔽、替换或恢复某组件后的变化，再与实际干预比较，并设置随机位置与等强度替换的对照。训练追溯则要挑选高影响、低影响及随机样本，通过小幅重新加权或部分重训练核对影响方向和大小；不能拿输入遮挡实验验证关于训练过程的论断。每类检验改变的对象不同，证据应与解释声称的作用对应。

模型生成的理由属于待检验输出。它可能概括了部分依据，也可能事后构造一个流畅理由。若模型在关键证据被替换后行为改变而理由保持不变，理由的忠实性就受到质疑。解释评价因此要连接模型行为，而不能只评价语言是否清晰。

\subsection{完整性与覆盖范围}
\label{sec:interpretability-quality-coverage}

忠实地解释一小部分行为，仍可能遗漏使用者真正关心的部分。这里需要区分数值分配的完整与行为范围的覆盖。对具有加性完整性要求的归因，可以计算
\[
 e_{\mathrm{sum}}
 =\left|f(\bm{x})-f(\bm{x}_0)-\sum_j a_j\right|,
\]
其中$a_j$为各项归因。积分梯度用这个残差检查数值积分，Shapley归因则用其声明的$v(\varnothing)$和$v(N)$替换两端分数。Grad-CAM并不承诺贡献和等于分数差，不应强行用这项标准判它失败。即使残差为零，也只表示所定义的数值差被分配完毕，未证明现实原因或全部内部通路已经解释清楚。

对整体代理或规则集，需要检验它能覆盖哪些输入和分支。可以预先划分语言、任务类型和重要错误情境，分别报告样本数、代理误差及规则能够给出解释的比例。若解释系统允许弃权，应同时报告有解释样本的质量与被覆盖比例，避免仅解释少数简单样本后宣称整体有效。测试集没有覆盖的区域，不能因为公式在该处也能计算就视为已经得到解释证据。

对机制解释，覆盖还包括相同现象是否在多个输入、不同表达和组合条件下成立。恢复某个位置可以支持一个局部机制假设；要概括一类行为，还需在新的输入配对上复测，并检查未被该机制解释的剩余效应。覆盖范围越大，需要承担的证据要求也越多。

\subsection{稳定性与一致性}
\label{sec:interpretability-quality-stability}

稳定性要检验的是在解释目标保持相同的条件下，结论会不会因计算随机性或无关变化而剧烈改变。固定模型、输入、基准和解释定义，重复邻域采样、背景采样或概念方向拟合，可以测量方法本身的波动。对数值归因可报告各项均值与标准差，对重要因素排序可比较秩相关，对前$k$项集合可计算交集大小除以并集大小。方向或特征编号可交换时，需要先匹配再比较，不能将编号不同当作语义不一致。

另一种试验是施加预期不改变解释目标的轻微输入变换，再检查解释是否相应保持。图像平移后的作用区域需要先映射回同一坐标，词语改写也需要先对应解释单位。即使语义和最终输出相同，模型也可能改用了另一组依据；因此，解释差异需要结合模型响应检查，不能仅凭输出未变就认定解释不稳定。如果模型本来就越过决策边界，解释改变更可能恰好忠实，不能一律惩罚。固定模型时的计算稳定性、换用新输入时的适用性和重新训练模型后的复现性，应分别报告。

还需检查解释是否对不该忽略的变化作出反应。随机化模型后解释完全不变，可能说明解释不依赖模型，而不是稳定性很好。反之，两套实现若计算同一函数，而方法声称只解释输入输出关系，则应检查是否得到相容结论。这些比较都需要先规定什么应保持、什么应改变，避免把“永远输出同一张图”当成高质量解释。

\subsection{可理解性}
\label{sec:interpretability-quality-comprehensibility}

可理解性关注目标读者能否正确读懂解释。规则数量、树深度、非零归因数量和概念名称熟悉程度，可以作为理解负担的线索，但只是替代指标。一条很短却含义模糊的规则，可能比稍长的明确计算更难理解。

直接检验可以让参与者只凭解释预测模型对新输入的反应，回答改变某个条件后模型会怎样判断，或指出规则适用的范围。记录答案正确率、完成时间和系统性误解，并控制任务难度、提供的事实和读者背景。例如，给出示例~\ref{ex:interpretability-global-local}的规则后，让读者判断$(0,1)$是否会被隔离，能检查其是否理解了加性分数与阈值；只问“这份解释清楚吗”，只能得到主观评价。

Doshi-Velez与Kim区分了真实应用中的人员评价、简化任务中的人员评价，以及不用人员参与的替代指标。这一区分说明，规则数少、归因稀疏等便于测量的指标，需要额外证据才能代表实际理解\scite{doshivelez2017rigorous}

在理解测试中，可以让使用者判断模型会不会拦截一组新邮件，以模型的实际行为核对答案。这与判断邮件本身是否应该拦截不同：前者检验人能否理解模型，后者还依赖其对任务的判断。纠正误拦和减少错误采纳的效果，留待下一小节放入使用任务中评价。只询问“是否喜欢这份解释”，则主要测量主观感受，不能替代这两种任务结果。

实验还应避免让解释组额外获得原本就缺失的关键信息，而对照组没有相同机会。否则改善可能来自信息量，而不是解释结构。可以给两组相同事实，仅改变组织方式，或单独比较附加事实与解释方法的贡献。对于不同专业背景的使用者，适合的粒度也可能不同；面向开发者有效的激活图，未必能帮助邮件收件人作出正确决定。

\subsection{使用有效性}
\label{sec:interpretability-quality-utility}

读懂模型不等于更好地使用模型。使用有效性要检验解释是否帮助发现错误、减少错误采纳或选择复核，并核算增加的时间与认知负担。补救行动执行后的效果由第~\ref{sec:interpretability-recourse-validation}节另行检验。因此，它需要把解释放回真实或足够接近真实的决策任务中。

可以比较无解释、所提解释与信息量相当的对照说明，让使用者处理同一难度的任务，记录最终决策损失、正确纠错比例、正确预测被无故否决的比例及用时。对同一人重复呈现同一题会带来记忆效应，设计时应随机分配题目或平衡呈现顺序。对存在模型错误的任务，还应单独报告错误建议被采纳的比例；总体采纳率提高可能只是过度信任。

例如，假设两个复核组各处理100封难度匹配的邮件，模型在各组都有20次错误。解释组纠正15次，但又把8个原本正确的判断改错；对照组纠正10次，只额外改错1次。最终错误数分别为$20-15+8=13$和$20-10+1=11$。这个教学例子说明，只报告解释组“纠错更多”会遗漏新的损害；需要比较最终结果，并用足够样本评估差异的不确定性。

五个维度应联合阅读：低代理误差提供忠实性证据，分组检查限定覆盖，重复试验测量稳定性，读者任务检验可理解性，实际决策对照检验使用收益。前一项通过不自动推出后一项；若因简化解释而放弃部分细节，应报告这种取舍落在哪些维度和输入范围中。

\section{决策的补救}
\label{sec:interpretability-recourse}

知道邮件为什么被误拦后，使用者还需要知道怎样恢复正常业务。第~\ref{sec:interpretability-local-counterfactual}节的反事实解释给出了能够改变模型判断的输入，但从“另一输入会被放行”到“使用者能够完成补救”，还缺少行动及其后果。\term{决策补救}（recourse）研究通过可行行动改变不利决策的办法；评价时还要检查这种改变是否符合实际任务目标。

本节围绕一份补救方案的形成与检验展开：先确定可以采取哪些行动、希望达到什么结果，再在行动约束下比较方案并选择，随后验证行动能否执行、后果是否符合预期、目标是否真正达到。这三个问题相互依赖。没有行动及其后果的描述，就无法判断约束下是否存在解；没有明确成本与目标，就无法比较哪种方案更好；求得最优解，也只说明它在所设模型中最优，仍需检验现实效果。

\subsection{补救行动的确定}
\label{sec:interpretability-recourse-target}

补救先要明确谁能够做什么。对于被误拦的邮件，发送方可以纠正错误的来源信息、补充可核验的证明，或在保留业务内容的前提下改变附件形式；收件人可能只能请求复核；系统运营方则能够检查提取错误或调整处理流程。不同参与者拥有不同权限，不能把一方能够完成的操作直接推荐给另一方。

这些行动可以从三个方向寻找。数据记录错误时，应纠正记录，使系统依据真实信息重新判断；记录正确但当前状态不满足要求时，可以改变可行动的条件，例如补充证明或完成必要培训；预测或处理流程存在错误时，可以提交申诉、请求人工复核。前两类通常能表示为行动后得到新输入，再由系统重新判断；人工复核还涉及决策流程的变化，需要说明由谁处理、依据什么以及何时反馈，不能等同于某个模型分数下降。

确定目标也不能只写“让结果翻转”。正常邮件的目标是保留真实业务内容并得到恰当处理，单纯删除付款要求虽然可能使评分下降，却可能把原业务请求一并删除。算法输出改善与任务得到补救应分别定义。若目标是进入人工复核，就应检查是否确实进入复核流程，而不是要求模型立即给出有利类别。

把候选行动写成具体方案时，应列出执行者、操作对象、所需资源、完成时间和预期变化。例如，“改善附件”不足以执行；“由发送方将附件转换为不含可执行内容且保留必要信息的格式，再重新提交”才说明了行动。在这里，解释用于提出候选方案，任务知识与实际权限决定它们是否可以实施。

为后续比较，设当前状态为$\bm{x}$，动作$a$实施后的状态由$T(\bm{x},a)$描述。$T$是行动后果模型：它说明执行动作后哪些字段会改变、改变多少，以及是否带来其他变化。固定预测规则$\hat y$时，候选结果为$\hat y(T(\bm{x},a))$；涉及人工复核等流程变更时，则需要另行描述流程结果。输入距离小不等于行动容易，直接改写记录也不等于改变了记录所描述的事实。

\subsection{补救行为的选择}
\label{sec:interpretability-recourse-action}

有了候选行动，还需要回答哪些能做、哪些值得做。不可改变的属性、操作权限、允许的变化方向和取值范围限定行动本身；费用、时间和业务内容的保留要求进一步限定行动后果。例如，附件必须保留合同信息、方案必须在截止时间前完成，都应作为筛选条件。必须满足的条件是硬约束，不能因为另一个方案分数更好就忽略；在可行方案之间偏好更少时间或更低费用，则属于优化目标。

设$A(\bm{x})$为满足这些行动约束的集合，$c(a;\bm{x})$为执行成本，$Y_{\mathrm{goal}}$为目标输出集合。对于行动后仍由固定规则判断的情形，最低成本补救可以写为
\begin{equation}
 \begin{aligned}
 \min_{a\in A(\bm{x})}\quad &c(a;\bm{x})\\
 \text{满足}\quad &\hat y(T(\bm{x},a))\in Y_{\mathrm{goal}}.
 \end{aligned}
 \label{eq:interpretability-recourse}
\end{equation}
$A$回答行动是否允许，$T$描述行动会造成什么变化，目标约束检查变化后是否得到所需结果，$c$用于比较达标方案。只有在目标约束下仍存在可行行动，才有相应的补救方案；其中成本最低者才是所设问题的最优解。这里的“最优”依赖行动范围、后果模型和成本定义，不表示对所有人、所有偏好都最好。

成本应反映实际负担，而不是直接相加不同单位的特征差。例如，时间、费用和操作次数可以分别报告；确需合为一个数时，应说明权重及其依据。若使用者对时间和费用有不同偏好，可以提供若干互不支配的方案，即不存在另一方案在所有所考虑的成本上都不更差、且至少一项更好，再由使用者选择。约束改变也可能改变最优方案，不能脱离可用资源推荐统一的“最小改动”。

Ustun等把可行动变量、变化方向、离散取值和成本写入线性分类器的补救问题，用整数规划搜索满足目标分类约束的行动，并评估哪些个体缺少可行补救。这类工作把输入层面的目标变化落实为可以检查和优化的行动条件\scite{ustun2019recourse}

实际求解取决于行动如何表示。候选方案较少时，可以逐项计算后果、筛除不满足约束者，再比较成本；连续与离散操作并存时，可以在适用的模型形式下使用数学规划或搜索。找到一个可行方案只证明存在这样的方案，有限搜索失败也不能自动证明无解。声称全局最优需要相应的搜索完备性或最优性证据，否则应说明得到的是当前搜索范围内的候选。

行动后果之间还可能存在依赖。培训增加技能，也占用工作时间，进而影响当期资源；优化器不能把这些变化独立设定。Karimi等提出在结构因果模型中计算干预后果，再在行动空间选择低成本干预。此时，改变一个变量还会按照结构方程影响其后继变量，区别于在输入表中独立修改多个字段\scite{karimi2021interventions}

\begin{example}[约束与成本决定补救方案]
\label{ex:interpretability-causal-recourse}
设一个教学模型仅在$s+r\geq3$时批准某项申请，其中$s$为技能分数，$r$为当期资源指标，当前状态为$(1,1)$。若只在输入表中把$s$改为2，就能越过边界；但假设培训强度$a\geq0$带来的实际变化是$s'=1+a$、$r'=1-a/2$，则总分为$2+a/2$，需要$a\geq2$。如果培训上限为$a\leq1$，仅靠培训就没有可行补救。

再假设可以获得强度为$b$的资源支持，$0\leq b\leq1$，使$r'=1-a/2+b$。允许组合两种行动，教学成本定义为$c(a,b)=a+3b$，则目标条件变成$a/2+b\geq1$。单独使用资源支持$(a,b)=(0,1)$可达标，成本为3；组合方案$(1,1/2)$也可达标，成本为2.5。

组合方案在上述假设下确实最优：任何可行解都满足$b\geq1-a/2$，因此$c(a,b)\geq3-a/2\geq2.5$，而$(1,1/2)$恰好达到该下界。若总预算进一步限定为2，则两种行动及其组合都无法在该预算下完成补救。这个结论依赖明确的行动范围和后果方程，而不是因为一次搜索没有找到方案。
\end{example}

若行动后果或未来模型不确定，可以要求方案在声明的一组情形下都有效，或在明确的概率模型下以足够高的概率达标。更严格的要求可能增加成本或使可行集合为空，应同时报告这种取舍。没有可行方案时，需要说明阻碍来自预算、权限、行动后果还是决策规则，并考虑复核或流程调整，不能要求使用者实施本就不可能的改动。

\subsection{补救效果的验证}
\label{sec:interpretability-recourse-validation}

优化得到的方案需要经过两层验证：在计算模型中确实满足约束和目标，以及在实际执行中产生了预期结果。前者核对求解是否正确，后者检验行动后果模型是否可信。第~\ref{sec:interpretability-quality-utility}节评价提供解释是否帮助人作出判断，本小节则评价执行补救行动后发生了什么，二者不能用同一项主观满意度代替。

\paragraph{方案与约束的复核}

对候选方案重新检查执行权限、不可变属性、变量取值、时间和预算，再按所声明的$T$生成新状态，通过完整的特征处理与决策流程重新计算结果。若求解时使用了代理模型或连续近似，必须回到实际系统复核，而不能只检查代理分数。成本最优性也应与可行性分开：可以用穷举、小规模精确求解或求解器提供的有效界限比较成本；无法证明最优时，应报告这一限制。

在示例~\ref{ex:interpretability-causal-recourse}中，代入$(a,b)=(1,1/2)$得到$(s',r')=(2,1)$，总分为3、成本为2.5，且满足两项行动上限。这验证了方案在教学模型中的可行性和达标性；前面的成本下界则另行证明最优性。若实际预算为2，同一方案即使能够达标，也不再可执行。

\paragraph{行动后果与实际收益的检验}

实际验证需要记录行动是否完成、实际花费多少、状态如何变化，以及最终决策是否达到目标。用实际新状态与$T$的预测作比较，才能发现后果模型是否高估了行动效果。只观察到技能与资源相关，通常不足以判断某个人培训后的变化；未观测因素和具体实施方式都会影响建议。可以结合已有干预证据、适当的对照研究或实际执行记录复核后果假设，并明确它们能支持的范围。

沿用同一例子，若实际每单位培训消耗的资源为0.8而非假设的0.5，那么执行$(1,1/2)$后状态为$(2,0.7)$，总分只有2.7，目标没有达到。求解器并未算错；失效来自行动后果模型。应据此修正后果估计并重新选择方案，而不能把纸面上的最优性当作实际有效性的证明。

评估一套补救方法时，应分别报告有可行建议的比例、建议被实际完成的比例、完成后达到目标的比例，以及实际成本和不利后果。每项比例都要说明分母；只在成功执行者中统计达标率，会漏掉没有方案或无法执行的人。若要声称建议本身改善了结果，还需与无建议或既有处理方式作适当对照，处理参与者自选择等差异，不能只比较执行前后的变化。

目标结果之外还要复核任务是否完成。邮件获得放行但合同内容被破坏，或模型评分降低却没有消除真实风险，都不能算完整补救。对于复核与申诉，应检查是否得到实质处理、原有错误是否纠正，以及耗时是否可接受。无论采用哪条路径，都要把决策变化与实际收益分别记录。

\paragraph{效果保持与失效处理}

行动往往需要时间。方案应注明依据的模型版本、后果假设和有效条件；执行周期较长时，需要检查模型更新、环境变化或后果估计误差是否会使它失效。可以在一组有依据的情形下重新计算达标与成本，或跟踪执行后的结果，但有限情形测试不能自动给出所有未来条件下的保证。

验证失败后应定位缺口：操作无法完成就修订行动集合，实际状态变化不符就修订后果模型，模型或流程变化则需要重新评估目标条件。三小节由此形成闭环：行动确定提供候选，约束优化选出方案，效果验证检验其依据并指导修订。某些个体最终仍可能没有可行补救，应明确说明限制和可用的复核途径。

\section{前沿研究}
\label{sec:interpretability-frontiers}

基础模型把解释对象扩展到高维内部表示和长行为过程。除了找到可读的模式，还需检验这些模式能否重复发现，以及局部解释能否连接成对完整行为的说明。

\subsection{从神经元到可检验的特征}

基础模型中的机制解释尝试从高维激活提取较可读的特征，再分析特征如何参与计算。直接给单个神经元命名往往不够，因为一个坐标可能对多种不相关情境响应，一种语义也可能分布在多个坐标。稀疏自编码器用一组数量更多的字典方向重构激活，并要求每个输入只使用少量方向。对激活$\bm h$、稀疏编码$\bm z=E(\bm h)$、字典矩阵$\bm D$和偏置$\bm b$，一个基本目标为
\begin{equation}
 \min_{E,\bm D,\bm b}
 \mathbb E_{\bm h}\!\left[
 \|\bm h-\bm D E(\bm h)-\bm b\|_2^2
 +\lambda\|E(\bm h)\|_1\right],
 \label{eq:interpretability-sae}
\end{equation}
并约束字典列的尺度，以免单纯放大字典、缩小编码来规避稀疏惩罚。重构项保留信息，稀疏项鼓励分开少数激活因素，但两者都没有直接保证人能够命名这些因素。

Bricken等在单层语言模型中研究了这种字典分解，并结合高激活样本与特征干预分析解释。相比按神经元检查，它允许语义方向偏离原始坐标轴；相应代价是字典宽度、稀疏程度、重构误差和特征拆分方式都需要选择。某个方向在若干样本上看似对应“危险内容”，不能因此把它当成完整的安全判别器\scite{bricken2023monosemanticity}

\subsection{特征解释的可复现性}

Song等研究稀疏自编码器在重复训练中的特征一致性，并提出在重构与稀疏性之外增加一致性评价。该结果把“发现了可读特征”推进到“发现是否能够复现”，但跨运行一致仍不单独证明特征具有现实因果意义或覆盖模型全部行为\scite{song2026featureconsistency}

一致性评价需要区分特征置换与实质分歧：两次运行中的同一语义可能对应不同编号，需要先匹配方向或激活模式；一个宽泛特征也可能在另一次训练中被拆成多个细粒度特征。可读性、重构质量、重复运行匹配和干预效果应共同报告，任何一个指标都不能单独代表机制解释已经完成。将解释用于第~\ref{sec:alignment-oversight}节的行为监督时，还需验证解释特征能否预测新的失败，而不是只复述已检查案例。

\subsection{跨步骤的行为依赖}

除了在重复训练中得到相近特征，长行为序列还需要解释跨步骤的依赖：局部合理的工具选择可能组合成失败。单次归因难以表达累积关系，需要结合机制干预与行为测试检验解释。

以邮件助手为例，误拦可能来自错误读取附件，也可能来自后续工具结果覆盖了原有证据。解释系统应明确时间顺序，区分读取到信息、在表示中保留信息和实际据此选择动作。若只对最终拒绝语句作归因，前面的错误状态可能已经不可见。面向这类任务，解释的目标将从一张静态贡献图扩展为可反复检验的行为依赖说明，同时保留干预本身的范围与限制。

\section{本章小结}
\label{sec:interpretability-summary}

对章首的误拦邮件，整体解释给出一组邮件上的规则与依据，单次解释则定位这一封邮件的判断。基于模型结构的解释交代计算规则，基于输入的解释比较字段或区域的作用，基于语义的解释检查这些依据对应的含义。三个方面相互联系，能够共同描述模型在所声明范围内怎样作出判断。

对一次判断，推理输入、推理计算过程和训练数据提供了三个互补的解释角度。相似案例、反事实、局部代理与特征归因共同研究固定模型对输入的响应；内部干预检验预测过程中的计算作用；训练数据追溯则研究学习经历如何影响当前结果。能够复算一份贡献，不代表已经找到内部机制；找到一条内部通路，也不代表已经说明现实因果关系。把多个单次结论提升为整体主张，还需明确汇总规则、异质性和样本覆盖。

解释质量由忠实性、完整性与覆盖、稳定性、可理解性和使用有效性共同约束，应以与论断对应的试验分别检验。即使这些检查通过，从“输入改变会改变输出”到“采取行动能够补救”，仍需可行动作与后果模型。解释清晰不代表模型正确，解释忠实也不代表依据公平。解释还可能展示训练样本、私人特征或内部记忆；为了让行为可理解而提供更多信息，也需要确定这些信息可以向谁公开。下一章将讨论数据使用、模型发布与信息披露中的隐私保护。
